\documentclass[preview]{standalone}

\usepackage{amsmath}
\usepackage{amssymb}
\usepackage{stellar}
\usepackage{definitions}

\begin{document}

\id{psicologia-cervello-comportamento}
\genpage

\section{Cervello e comportamento}

\begin{snippet}{biologia-comportamento-metodi}
    I neuroscienziati studiano il rapporto tra cervello e comportamento a
    molteplici livelli. Alcuni metodi osservano casi di lesione cerebrale,
    altri producono interferenze temporanee e controllate, altri ancora
    registrano l'attività cerebrale durante compiti cognitivi o comportamentali.
\end{snippet}

\begin{snippetexample}{phineas-gage-example}{Phineas Gage}
    Nel 1848 il caposquadra ferroviario Phineas Gage subì un grave incidente:
    una barra di metallo gli attraversò la testa in seguito a un'esplosione.
    Dopo l'incidente conservò molte abilità motorie e linguistiche, ma mostrò
    cambiamenti marcati nella personalità e nel comportamento razionale. Il
    caso suggerì l'esistenza di relazioni specifiche tra danno cerebrale e
    funzioni psicologiche.
\end{snippetexample}

\begin{snippetdefinition}{area-broca-definition}{Area di Broca}
    L'\emph{area di Broca} è una regione cerebrale frontale sinistra cruciale
    per la produzione del linguaggio verbale.
\end{snippetdefinition}

\begin{snippetdefinition}{stimolazione-magnetica-transcranica-definition}{Stimolazione magnetica transcranica}
    La \emph{stimolazione magnetica transcranica} è una tecnica che utilizza
    impulsi magnetici per interferire temporaneamente con l'attività di aree
    cerebrali specifiche, producendo effetti reversibili simili a lesioni
    funzionali controllate.
\end{snippetdefinition}

\begin{snippetdefinition}{elettroencefalogramma-definition}{Elettroencefalogramma}
    L'\emph{elettroencefalogramma}, o \emph{EEG}, è una registrazione
    dell'attività elettrica cerebrale ottenuta tramite elettrodi posizionati
    sulla superficie del cuoio capelluto.
\end{snippetdefinition}

\begin{snippetdefinition}{tomografia-emissione-positroni-definition}{Tomografia a emissione di positroni}
    La \emph{tomografia a emissione di positroni}, o \emph{PET}, è una tecnica
    di neuroimaging che registra la radioattività delle cellule per ottenere
    immagini dell'attività cerebrale durante compiti cognitivi e
    comportamentali.
\end{snippetdefinition}

\begin{snippetdefinition}{risonanza-magnetica-definition}{Risonanza magnetica}
    L'\emph{imaging con risonanza magnetica}, o \emph{MRI}, è una tecnica di
    neuroimaging che usa campi magnetici e onde radio per ottenere immagini
    dettagliate della struttura anatomica del cervello.
\end{snippetdefinition}

\begin{snippetdefinition}{encefalo-definition}{Encefalo}
    L'\emph{encefalo} è la componente principale del sistema nervoso centrale
    contenuta nella scatola cranica. Comprende strutture interconnesse che
    regolano processi vitali, comportamento emotivo, motivazione e funzioni
    cognitive superiori.
\end{snippetdefinition}

\begin{snippet}{strutture-principali-encefalo}
    L'encefalo comprende tre grandi strutture funzionali: il
    \emph{tronco encefalico}, che regola processi vitali come il battito
    cardiaco; il \emph{sistema limbico}, coinvolto nel comportamento emotivo e
    negli impulsi motivazionali; il \emph{telencefalo}, coinvolto nelle funzioni
    cognitive ed emotive di livello superiore.
\end{snippet}

\includesnpt[width=60\%,src=/snippet/static-psychology/major-brain-structures.jpg]{centered-img}

\begin{snippetdefinition}{corpo-calloso-definition}{Corpo calloso}
    Il \emph{corpo calloso} è un fascio di fibre nervose che trasferisce
    informazioni tra i due emisferi cerebrali in entrambe le direzioni.
\end{snippetdefinition}

\begin{snippetdefinition}{pazienti-split-brain-definition}{Pazienti split-brain}
    I \emph{pazienti split-brain} sono pazienti nei quali il corpo calloso è
    stato reciso chirurgicamente, di solito come trattamento di forme gravi di
    epilessia.
\end{snippetdefinition}

\begin{snippet}{vie-visive-split-brain}
    Nei pazienti split-brain, la recisione del corpo calloso impedisce lo
    scambio diretto di informazioni tra i due emisferi. Poiché nella maggior
    parte delle persone il linguaggio è prevalentemente controllato
    dall'emisfero sinistro, gli stimoli elaborati solo dall'emisfero destro
    possono influenzare l'azione manuale senza poter essere nominati
    verbalmente.
\end{snippet}

\includesnpt[width=65\%,src=/snippet/static-psychology/split-brain-visual-pathways.jpg]{centered-img}

\begin{snippetexample}{paziente-split-brain-example}{Paziente split-brain}
    Se un oggetto appare brevemente nel campo visivo sinistro di un paziente
    split-brain, l'informazione raggiunge l'emisfero destro. Il paziente può
    riconoscere l'oggetto con la mano sinistra, controllata prevalentemente
    dall'emisfero destro, ma non riesce a nominarlo verbalmente perché il
    linguaggio dipende soprattutto dall'emisfero sinistro. Se invece deve
    nominare un oggetto elaborato dall'emisfero sinistro, la risposta verbale è
    possibile.
\end{snippetexample}

\begin{snippetdefinition}{lateralizzazione-definition}{Lateralizzazione}
    Una funzione è \emph{lateralizzata} quando, per la sua realizzazione, un
    emisfero cerebrale svolge un ruolo primario rispetto all'altro.
\end{snippetdefinition}

\begin{snippet}{lateralizzazione-non-dicotomica}
    La lateralizzazione non implica una separazione rigida tra emisfero sinistro
    ed emisfero destro. L'emisfero sinistro ha un ruolo predominante in molte
    funzioni linguistiche, ma l'esperienza psicologica dipende dall'azione
    integrata dei due emisferi.
\end{snippet}

\begin{snippetdefinition}{plasticita-cerebrale-definition}{Plasticità cerebrale}
    La \emph{plasticità cerebrale} è la capacità del cervello di modificare le
    proprie prestazioni e la propria organizzazione nel corso del tempo, anche
    attraverso la formazione di nuove sinapsi o la modificazione di sinapsi già
    esistenti.
\end{snippetdefinition}

\begin{snippetdefinition}{neurogenesi-definition}{Neurogenesi}
    La \emph{neurogenesi} è la produzione di nuove cellule cerebrali, cioè di
    nuovi neuroni.
\end{snippetdefinition}

\end{document}
